\section{Базовые модели, промежуточные эксперименты и выбор финальной модели}

\subsection{Роль baseline-моделей}

Построение рекомендательной системы по истории культур требует не только обучения основной модели, но и сравнения её с более простыми ориентирами. Без такого сравнения невозможно понять, действительно ли итоговое качество связано с использованием более содержательной модели, а не достигается за счёт простых частотных закономерностей или инерционных эффектов последовательности.

В данной работе baseline-модели выполняют несколько функций. Во-первых, они задают нижнюю границу качества, с которой можно сопоставлять более сложные методы. Во-вторых, они позволяют проверить, насколько полезными оказываются история поля и дополнительные признаки по сравнению с простыми правилами. В-третьих, baseline-блок помогает обосновать выбор финальной модели не декларативно, а на основе сравнительного анализа.

По этой причине в работе рассматриваются:

\begin{itemize}
\tightlist
\item
  простые baseline-подходы;
\item
  вероятностный baseline Markov-1;
\item
  baseline-модель CatBoost;
\item
  промежуточные экспериментальные ветки.
\end{itemize}

Такой набор делает выбор итоговой модели методически прозрачным и воспроизводимым.

\subsection{Простые baseline-подходы}

В качестве простых ориентиров используются два базовых подхода: \texttt{most\_frequent} и \texttt{last\_crop}.

\subsubsection{\texorpdfstring{Baseline \texttt{most\_frequent}}{Baseline most\_frequent}}

Модель \texttt{most\_frequent} всегда предсказывает наиболее частый класс в обучающей выборке. Этот подход не использует ни историю конкретного поля, ни региональный контекст, ни пространственные признаки. Его практическая ценность невелика, однако он важен как минимальный ориентир: если более сложная модель не превосходит такой baseline, то её использование трудно считать обоснованным.

\subsubsection{\texorpdfstring{Baseline \texttt{last\_crop}}{Baseline last\_crop}}

Модель \texttt{last\_crop} в качестве следующей культуры выбирает последнюю наблюдаемую культуру в истории поля. В отличие от \texttt{most\_frequent}, этот baseline уже учитывает локальную информацию о поле, но делает это в крайне упрощённой форме. Он отражает инерционный сценарий, в котором предполагается, что следующая культура совпадает с ближайшим предыдущим состоянием.

Сравнение с этим baseline особенно важно для рассматриваемой задачи, поскольку в последовательностях культур действительно могут встречаться повторы, а значит, правило последней культуры способно давать нетривиальный результат.

\subsection{Markov-1 как вероятностный baseline}

Следующий baseline-подход строится на одношаговой вероятностной модели переходов между культурами. В этой постановке вероятность следующей культуры зависит только от последней культуры в истории:

\begin{center}\(
P(\mathrm{target}\mid\mathrm{history}_3)
\)\end{center}

Такой подход уже не является чисто эвристическим: он использует статистику переходов, наблюдаемую в обучающей выборке. При этом он по-прежнему остаётся ограниченным, поскольку игнорирует более длинную историю поля, а также дополнительные признаки региона и участка.

Markov-1 играет важную роль в работе. С одной стороны, это сильнее, чем наивные baseline-стратегии. С другой стороны, он остаётся достаточно простым и интерпретируемым, чтобы служить честным ориентиром для сравнения с CatBoost. В дальнейшем именно Markov-1 используется как основной baseline для recommendation-сравнения по метрикам \texttt{Accuracy@1}, \texttt{Hit@3} и \texttt{Hit@5}.

\subsection{Baseline CatBoost-модель}

Центральным предиктивным компонентом системы является baseline-версия модели CatBoost. Выбор этой модели связан с характером задачи: данные имеют табличную структуру и содержат как категориальные, так и числовые признаки. История культур и региональные коды естественным образом выступают как категориальные признаки, а площадь поля и координаты --- как числовой контекст.

Baseline CatBoost использует компактный набор признаков:

\begin{itemize}
\tightlist
\item
  \texttt{history\_1}, \texttt{history\_2}, \texttt{history\_3};
\item
  \texttt{STATEFIPS}, \texttt{CNTYFIPS};
\item
  \texttt{CSBACRES};
\item
  \texttt{INSIDE\_X}, \texttt{INSIDE\_Y}.
\end{itemize}

Такой набор отражает основную гипотезу работы: следующая культура определяется не только последней наблюдаемой культурой, но и более длинной историей, а также региональным и пространственным контекстом. При этом baseline-модель остаётся достаточно компактной, чтобы быть воспроизводимой и удобной для последующего анализа.

Финальная baseline-модель на test показывает:

\begin{itemize}
\tightlist
\item
  \texttt{Accuracy@1\ =\ 0.69901};
\item
  \texttt{Macro\ F1\ =\ 0.56849};
\item
  \texttt{Weighted\ F1\ =\ 0.68963}.
\end{itemize}

Эти значения показывают, что CatBoost уже на уровне базовой конфигурации существенно превосходит более простые ориентиры. Тем самым baseline-модель становится кандидатом на роль финального предиктивного ядра системы.

\subsection{Эксперимент с дополнительными признаками}

После построения baseline-модели был проведён отдельный эксперимент с расширением признакового пространства. Цель этого этапа состояла в том, чтобы проверить, может ли модель выиграть от признаков, описывающих структуру истории культур более явно, чем исходные три года истории. Такой эксперимент был реализован в отдельной ветке \texttt{05\_catboost\_improved.ipynb}.

В эксперименте использовались дополнительные derived-признаки, в том числе:

\begin{itemize}
\tightlist
\item
  число уникальных культур в истории;
\item
  наличие \texttt{fallow} в истории;
\item
  наличие \texttt{legumes} в истории;
\item
  признаки повторяемости соседних лет;
\item
  индикатор полного повтора истории.
\end{itemize}

Содержательно эта ветка проверяла гипотезу о том, что явное кодирование некоторых паттернов истории может дополнительно усилить модель. Однако по итогам проекта данное направление не было выбрано как основное. Причина состоит не в том, что эксперимент оказался бесполезным, а в том, что он не дал достаточных оснований для смены финального проектного контракта. В результате improved-ветка рассматривается как исследовательский эксперимент, полезный для понимания возможных направлений улучшения, но не как финальная модель диплома.

\subsection{Эксперимент с CatBoost и transition graph}

Следующий промежуточный эксперимент был связан с попыткой усилить ранжирование CatBoost за счёт графового сигнала переходов между культурами. Идея состояла в том, чтобы объединить модельные вероятности и вероятности, полученные из transition graph, по формуле

\begin{center}\(
\mathrm{score}=\alpha\,p_{\mathrm{catboost}}+\beta\,p_{\mathrm{graph}}
\)\end{center}

где \texttt{p\_catboost} --- вероятность, рассчитанная моделью, а \texttt{p\_graph} --- графовый сигнал переходов. В эксперименте проверялись несколько сочетаний весов: \texttt{(0.9,\ 0.1)}, \texttt{(0.8,\ 0.2)} и \texttt{(0.7,\ 0.3)}. Лучшим вариантом оказалось соотношение \texttt{alpha\ =\ 0.9}, \texttt{beta\ =\ 0.1}.

Однако даже в лучшей конфигурации гибридный вариант на test показал:

\begin{itemize}
\tightlist
\item
  \texttt{Accuracy\ =\ 0.69871};
\item
  \texttt{Macro\ F1\ =\ 0.56439};
\item
  \texttt{Weighted\ F1\ =\ 0.68873};
\item
  \texttt{Hit@3\ =\ 0.94956}.
\end{itemize}

Эти значения оказались немного хуже, чем у чистого CatBoost. Следовательно, простой graph reranking не был принят как финальная стратегия. В то же время сам эксперимент имеет важное методическое значение. Он показывает, что графовый сигнал можно интегрировать в модельный pipeline, но в текущей конфигурации его полезнее использовать не как механизм замены ранжирования, а как источник explainability-контекста для shortlist-кандидатов. Именно этот отрицательный результат позже становится важным основанием для архитектурного разведения предиктивного ядра и knowledge graph explainability layer.

\begin{center}
\small
\begin{longtable}{p{0.26\textwidth} p{0.31\textwidth} p{0.33\textwidth}}
\caption{Промежуточные эксперименты и выбор финальной модели}
\label{tab:intermediate-experiments}\\
\toprule
Эксперимент & Что проверялось & Итог \\
\midrule
\endfirsthead

\toprule
Эксперимент & Что проверялось & Итог \\
\midrule
\endhead

\bottomrule
\endlastfoot

Baseline CatBoost &
Компактная модель на признаках истории, региона, площади и координат поля &
\textbf{Выбран.} Зафиксирован как активная финальная модель в \path{artifacts/model_registry.json}; на test: Accuracy = 0.699006, Macro F1 = 0.568485, Weighted F1 = 0.689626 \\

CatBoost + extra features &
Расширение признакового пространства derived-признаками истории: число уникальных культур, наличие \texttt{fallow} и \texttt{legumes}, признаки повторяемости &
\textbf{Не выбран.} На test: Accuracy = 0.699235, Macro F1 = 0.569473, Weighted F1 = 0.689991. Несмотря на близкое качество, ветка сохранена как экспериментальная \texttt{improved} и не заменяет baseline в финальном проектном контракте \\

CatBoost + transition graph &
Гибридное ранжирование по формуле $\mathrm{score}=\alpha p_{\mathrm{catboost}}+\beta p_{\mathrm{graph}}$ с лучшей конфигурацией $\alpha=0.9$, $\beta=0.1$ &
\textbf{Не выбран.} Лучший hybrid-вариант на test немного уступил чистому CatBoost: Accuracy = 0.698714, Macro F1 = 0.564387, Weighted F1 = 0.688729; graph-сигнал оставлен как explainability layer \\

\end{longtable}
\footnotesize
Источники: \path{artifacts/model_registry.json},
\path{artifacts/catboost/improved/meta.json},
\path{notebooks/05_catboost_improved.ipynb},
\path{artifacts/results/classification_quality/catboost_baseline_metrics.csv},
\path{artifacts/results/graph_hybrid/hybrid_metrics.csv}.
\normalsize
\end{center}


\begin{center}
\scriptsize
\begin{longtable}{p{0.24\textwidth} p{0.31\textwidth} r r r}
\caption{Сравнение промежуточных экспериментальных конфигураций на тестовой выборке}
\label{tab:intermediate-experiments-metrics}\\
\toprule
Эксперимент & Что проверялось & Accuracy & Macro F1 & Weighted F1 \\
\midrule
\endfirsthead

\toprule
Эксперимент & Что проверялось & Accuracy & Macro F1 & Weighted F1 \\
\midrule
\endhead

\bottomrule
\endlastfoot

Baseline CatBoost &
История, регион, площадь и координаты поля &
0.699006 & 0.568485 & 0.689626 \\

CatBoost + extra features &
Расширение признакового пространства: число уникальных культур, наличие \texttt{fallow} и \texttt{legumes}, признаки
повторяемости &
0.699235 & 0.569473 & 0.689991 \\

CatBoost + transition graph &
Гибридное ранжирование по формуле
\(\mathrm{score}(c)=\alpha p_{\mathrm{CatBoost}}(c\mid x)+\beta p_{\mathrm{graph}}(c\mid c_{t-1})\)
при лучшей конфигурации \(\alpha=0.9\), \(\beta=0.1\) &
0.698714 & 0.564387 & 0.688729 \\

\end{longtable}
\end{center}


\subsection{Сравнение подходов и выбор финальной модели}

Сравнение основных baseline-подходов показывает, что CatBoost baseline заметно превосходит простые и вероятностные ориентиры по ключевым классификационным метрикам. Для validation и test были получены результаты, приведённые в таблице~\ref{tab:baseline-catboost-comparison}.

\begin{center}
\small
\begin{longtable}{p{0.25\textwidth} p{0.18\textwidth} r r r}
\caption{Сравнение baseline-подходов и финальной CatBoost-модели}
\label{tab:baseline-catboost-comparison}\\
\toprule
Модель & Выборка & Accuracy & Macro F1 & Weighted F1 \\
\midrule
\endfirsthead

\toprule
Модель & Выборка & Accuracy & Macro F1 & Weighted F1 \\
\midrule
\endhead

\bottomrule
\endlastfoot

\texttt{most\_frequent} & validation & 0.305306 & 0.051977 & 0.142820 \\
\texttt{last\_crop} & validation & 0.354267 & 0.172492 & 0.361115 \\
\texttt{markov\_1} & validation & 0.557140 & 0.411203 & 0.539053 \\
\texttt{catboost} & validation & 0.699281 & 0.568771 & 0.689966 \\

\texttt{most\_frequent} & test & 0.305068 & 0.051946 & 0.142623 \\
\texttt{last\_crop} & test & 0.354321 & 0.172575 & 0.361135 \\
\texttt{markov\_1} & test & 0.557219 & 0.411371 & 0.539027 \\
\texttt{catboost} & test & 0.699006 & 0.568485 & 0.689626 \\

\end{longtable}
\footnotesize
Источник значений: \path{artifacts/results/classification_quality/baseline_catboost_comparison_metrics.csv}.
\normalsize
\end{center}


Из этой таблицы следует несколько выводов.

Во-первых, простые baseline-стратегии действительно полезны как ориентиры, но недостаточны как основа рекомендательной системы. Во-вторых, Markov-1 оказывается заметно сильнее наивных эвристик, что подтверждает полезность переходной статистики между культурами. В-третьих, baseline CatBoost устойчиво превосходит все baseline-подходы, что подтверждает полезность трёхлетней истории и контекстных признаков. Наконец, промежуточные эксперименты с дополнительными признаками и с графовым reranking не дали достаточных оснований для отказа от baseline-модели как финального варианта.

Следовательно, именно baseline CatBoost обоснованно выбирается в качестве финального предиктивного ядра системы. Это решение согласуется и с проектным контрактом репозитория, где baseline зафиксирован как основная модель, а improved и graph-based экспериментальные ветки рассматриваются как non-final направления исследования.

\subsection{Выводы по главе}

В данной главе были рассмотрены baseline-подходы, промежуточные эксперименты и логика выбора финальной модели. Полученные результаты показывают, что:

\begin{itemize}
\tightlist
\item
  простые baseline-стратегии полезны как нижняя граница качества, но недостаточны как основа системы;
\item
  Markov-1 является сильным интерпретируемым ориентиром, однако уступает модели, использующей более длинную историю и контекстные признаки;
\item
  baseline CatBoost демонстрирует наилучшее качество и поэтому выбран как финальная модель;
\item
  эксперимент с дополнительными признаками и эксперимент с \texttt{CatBoost\ +\ transition\ graph} важны как исследовательские ветки, но не меняют итоговый выбор;
\item
  графовый сигнал оказался полезнее как explainability-компонент, чем как механизм reranking.
\end{itemize}

Таким образом, к концу данного этапа в работе зафиксировано финальное предиктивное ядро, на основе которого далее строятся recommendation layer, confidence-анализ и explainability-слой.
